\documentclass{article}
\usepackage{amsmath, amssymb, amsthm}
\usepackage{hyperref}
\usepackage{geometry}
\geometry{margin=1in}

\title{Erd\H{o}s Problem \#87}
\date{Last updated: 2025-08-31}
\author{Source: erdosproblems.com}

\begin{document}
\maketitle

\noindent\textbf{Status:} OPEN \quad \textbf{No prize}

\noindent\textbf{Tags:} graph theory, ramsey theory

\medskip

\noindent\textbf{Problem Statement:}

Let $\epsilon &#62;0$. Is it true that, if $k$ is sufficiently large, then\[R(G)&#62;(1-\epsilon)^kR(k)\]for every graph $G$ with chromatic number $\chi(G)=k$? Even stronger, is there some $c&#62;0$ such that, for all large $k$, $R(G)&#62;cR(k)$ for every graph $G$ with chromatic number $\chi(G)=k$?




Erd\H{o}s originally conjectured that $R(G)\geq R(k)$, which is trivial for $k=3$, but fails already for $k=4$, as Faudree and McKay \cite{FaMc93} showed that $R(W)=17$ for the pentagonal wheel $W$. Since $R(k)\leq 4^k$ this is trivial for $\epsilon\geq 3/4$. Yuval Wigderson points out that $R(G)\gg 2^{k/2}$ for any $G$ with chromatic number $k$ (via a random colouring), which asymptotically matches the best-known lower bounds for $R(k)$. This problem is #12 and #13 in Ramsey Theory in the graphs problem collection.




References

[FaMc93] Faudree, R. J. and McKay, B., A conjecture of Erd\H{o}s and the Ramsey number $r(W_6)$ . J. Combinatorial Math. and Combinatorial Computing (1993), 23-31.

\noindent\textbf{Related OEIS sequences:} A059442, possible


\bigskip
\noindent\small{Source: \url{https://www.erdosproblems.com/87}}
\end{document}
